\documentclass{article}
\newcommand{\dom}{\mathrm{dom}}
\usepackage{lineno}
\linenumbers

% if you need to pass options to natbib, use, e.g.:
%     \PassOptionsToPackage{numbers, compress}{natbib}
% before loading neurips_2020

% ready for submission
% \usepackage{neurips_2020}

% to compile a preprint version, e.g., for submission to arXiv, add add the
% [preprint] option:
%     \usepackage[preprint]{neurips_2020}

% to compile a camera-ready version, add the [final] option, e.g.:
%     \usepackage[final]{neurips_2020}

% to avoid loading the natbib package, add option nonatbib:
%     \usepackage[nonatbib]{neurips_2020}

\usepackage[utf8]{inputenc} % allow utf-8 input
\usepackage[T1]{fontenc}    % use 8-bit T1 fonts
\usepackage{hyperref}       % hyperlinks
\usepackage{url}            % simple URL typesetting
\usepackage{booktabs}       % professional-quality tables
\usepackage{amsfonts}       % blackboard math symbols
\usepackage{nicefrac}       % compact symbols for 1/2, etc.
\usepackage{microtype}      % microtypography
\usepackage{amssymb,amsmath,amsthm}
\usepackage[framemethod=tikz]{mdframed}
\usepackage{verbatim}
\usepackage{bbm}

%\usepackage{mystyle}
\usepackage{mathtools,booktabs,tabu,graphicx,hyperref,epstopdf,url,rotating}
%\usepackage{fullpage}
\usepackage{algorithm}
\usepackage{cite}
\usepackage{algorithmic}
\usepackage{fancyref}
\usepackage[framemethod=tikz]{mdframed}

\newtheorem{theorem}{Theorem}[section]
\newtheorem{corollary}{Corollary}[theorem]	
\newtheorem{condition}{Condition}[theorem]	
\newtheorem{proposition}{Proposition}[theorem]
\newtheorem{definition}[theorem]{Definition}
\newtheorem{example}[theorem]{Example}
\newtheorem{assumption}[theorem]{Assumption}
\newtheorem{lemma}[theorem]{Lemma}
\newtheorem{remark}{Remark}[theorem]

\newcommand{\tstep}{{\Delta t}}

\renewcommand{\vec}[1]{{\mathbf{#1}}} 
\newcommand{\Oo}{\mathcal{O}}
\newcommand{\Nn}{\mathcal{N}}
\newcommand{\Id}{\mathrm{Id}}
\newcommand{\todolist}[1]{\textcolor{red}{(TODO: #1)}}
\newcommand{\eqdef}{\coloneq}
\newcommand{\Prox}{\mathrm{Prox}}

\newcommand{\pa}[1]{\pp{#1}}
\newcommand{\PP}{\mathbb{P}}
\newcommand{\EE}{\mathbb{E}}
\newcommand{\RR}{\mathbb{R}}

\newcommand{\clip}{L}
\newcommand{\clower}{K}
\newcommand{\cinfty}{B}
\newcommand{\ccc}{\clower}

%\usepackage{graphicx}
%\usepackage{subfigure}
\usepackage{subcaption}
\usepackage{geometry}
\geometry{
 a4paper,
 total={170mm,257mm},
 left=20mm,
 top=20mm,
 }

\DeclareMathOperator*{\diag}{diag}
\DeclareMathOperator*{\argmax}{arg\,max\,}
\DeclareMathOperator*{\argmin}{arg\,min\,}
\DeclareMathOperator*{\tr}{tr}


\title{Convergence of Numerical Discretizations of SDEs}
% The \author macro works with any number of authors. There are two commands
% used to separate the names and addresses of multiple authors: \And and \AND.
%
% Using \And between authors leaves it to LaTeX to determine where to break the
% lines. Using \AND forces a line break at that point. So, if LaTeX puts 3 of 4
% authors names on the first line, and the last on the second line, try using
% \AND instead of \And before the third author name.

\author{%
 Ivan Cheltsov
}

\newcommand{\OO}{\mathbb{O}}

\newcommand{\commentout}[1]{}


\begin{document}

\maketitle

\subsection{Almost Sure Convergence of Time Averages}

Consider a stochastic process $(x_t)_{t \geq 0}$ generated by an elliptic or hypoelliptic operator $\mathcal{L}$, which has a unique invariant probability measure $\pi$ with positive density relative to the Lebesgue measure. For any observable $\varphi \in L^1(\pi)$, the time average
\[
\hat{\varphi}_t = \frac{1}{t} \int_0^t \varphi(x_s) \, ds
\]
converges almost surely, for any starting point $x_0 \in \mathcal{X}$, to the true average
\[
\lim_{t \to +\infty} \hat{\varphi}_t = \int_{\mathcal{X}} \varphi(x) \, \pi(dx),
\]
as shown in \cite{kliemann1987}.  
This result not only guarantees that time-averaged estimates are accurate, but also proves the uniqueness of the invariant measure.
\subsection{Asymptotic Variance and Comparison to i.i.d. Sampling}

In addition to convergence, we want to understand how quickly the variance of $\hat{\varphi}_t$ decays. If $x_0 \sim \pi$, the asymptotic variance is given by

\[
\varsigma_\varphi^2 = \lim_{t \to +\infty} t \, \mathrm{Var}_\pi(\hat{\varphi}_t),
\]
which can be written as \cite{Lelievre2016-hx}
\[
\varsigma_\varphi^2 = -2 \int_{\mathcal{X}} (\mathcal{L}^{-1} \Pi \varphi)(x) \, (\Pi \varphi)(x) \, \pi(dx),
\]
where $\Pi \varphi = \varphi - \int \varphi \, d\pi$ represents the centered version of $\varphi$. This variance $\varsigma_\varphi^2$ captures how long the process takes to `forget' its starting point. It can be compared to the variance from averaging i.i.d. samples $(x_n) \sim \pi$, where the variance is simply
\[
\varsigma_{\varphi, \text{iid}}^2 = \int_{\mathcal{X}} (\Pi \varphi)^2 \, d\pi.
\]
Thus, the relation
\[
\varsigma_\varphi^2 = \theta_{\text{corr}, \varphi} \, \varsigma_{\varphi, \text{iid}}^2
\]
defines the correlation time $\theta_{\text{corr}, \varphi}$, which indicates how much slower the convergence is compared to independent sampling. Thus, to get an estimator of the same quality as $N$ i.i.d. samples, the integration time needs to be $t \approx N \theta_{\text{corr}, \varphi}$. When $\mathcal{L}$ is self-adjoint, the correlation time is controlled by the spectral gap $\lambda_1 > 0$ of $-\mathcal{L}$, with the bound
\[
\theta_{\text{corr}, \varphi} \leq \frac{1}{\lambda_1},
\]
and equality holds when $\varphi$ is the eigenfunction corresponding to $\lambda_1$.

\section{Convergence of SDE Discretizations}

\subsection{Introduction}

Consider a stochastic process $(x_t)_{t \geq 0} \subset \mathbb{R}^d$ defined as the solution to the Itô stochastic differential equation
\[
dx_t = b(x_t)\,dt + \sigma(x_t)\,dW_t,
\]
with initial condition $x_0 \in \mathbb{R}^d$, drift $b: \mathbb{R}^d \to \mathbb{R}^d$, diffusion $\sigma: \mathbb{R}^d \to \mathbb{R}^{d \times m}$ and $W_t \in \mathbb{R}^m$ a standard $m$-dimensional Brownian motion. In practice, such processes are rarely solvable in closed form and must be approximated via numerical discretization schemes. For example, a commonly used time-stepping scheme is the Euler–Maruyama method. For a fixed timestep $\Delta t > 0$, the scheme is given by
\[
x^{n+1} = x^n + b(x^n)\Delta t + \sigma(x^n)\sqrt{\Delta t}\, G^n,
\]
where $G^n \sim \mathcal{N}(0, I_m)$ are independent standard Gaussian vectors. Such a recursion defines a Markov chain with transition operator
\begin{equation} \label{discrete_transition}
    P_{\Delta t} \varphi(x) := \mathbb{E}[\varphi(x^{n+1}) \mid x^n = x],
\end{equation}
which can be viewed as a discrete analogue of the semigroup $(e^{t\mathcal{L}})_{t \geq 0}$ associated with the generator $\mathcal{L}$ of the continuous process.

\subsection{Weak and Strong Errors}
When evaluating the accuracy of an SDE discretization scheme, two primary notions of convergence are considered: weak and strong error.
Weak convergence is concerned with the convergence of expectations of functionals of the solution, which is often sufficient in applications such as statistical estimation or sampling. In contrast, strong convergence focuses on path-wise accuracy, measuring how closely the trajectories of the numerical approximation follow those of the exact solution. The distinction between these two types of convergence plays a critical role in choosing and analyzing discretization methods.

\subsubsection{Weak Convergence}

Let $\varphi \in C^\infty_c(\mathbb{R}^d)$ be a smooth test function with compact support. Then, under suitable regularity conditions (e.g., Lipschitz continuity of $b$ and $\sigma$), there exists $\alpha > 0$ such that, for any fixed terminal time $T > 0$, there are constants $C > 0$ and $\Delta t^* > 0$ satisfying
\[
\sup_{0 \leq n \leq T/\Delta t} \left| \mathbb{E}[\varphi(x^n)] - \mathbb{E}[\varphi(x_{n\Delta t})] \right| \leq C \Delta t^\alpha,
\quad \text{for all } \Delta t \in (0, \Delta t^*].
\]
This result quantifies the weak convergence of the numerical method, i.e., the convergence of expected values. The order $\alpha$ is typically determined via expansions such as
\[
P_{\Delta t} = e^{\Delta t \mathcal{L}} + \mathcal{O}(\Delta t^{\alpha+1}),
\]
and reflects the accuracy with which the discretization preserves average properties of the true process.

\subsubsection{Strong Convergence}

Stronger notions of convergence compare sample paths directly. The discretization satisfies a strong convergence estimate if there exists $\alpha > 0$ such that, for all $T > 0$, there are constants $C > 0$ and $\Delta t^* > 0$ with
\[
\sup_{0 \leq n \leq T/\Delta t} \left( \mathbb{E} \left| x^n - x_{n\Delta t} \right|^p \right)^{1/p} \leq C \Delta t^\alpha,
\quad \text{for all } \Delta t \in (0, \Delta t^*].
\]
Such estimates are often derived using stochastic versions of the Gr\"{o}nwall inequality. The strong error controls the pathwise approximation of the numerical scheme to the true solution.
%\subsubsection{Comparison of Weak and Strong Convergence}
%While both notions of convergence are useful, they serve different purposes:
%\begin{itemize}
%    \item Weak convergence is relevant for approximating expectations of functionals of the process.
%    \item Strong convergence is essential for path-dependent problems or when sample path accuracy is required.
%\end{itemize}
It is important to note that the constants in these error bounds are generally not uniform in time. Therefore, for long-term integration or ergodic sampling, further techniques are needed to control the accumulation of error over time.

\begin{example}
    For SDEs with globally Lipschitz drift and diffusion coefficients, the Euler–Maruyama method is known to achieve weak order $\alpha = 1$ and strong order $\alpha = 1/2$. 
\end{example}

\subsection{Convergence of Time Averages}
Weak and strong error estimates provide insights into the transient behavior of a discretization but are not sufficient for analyzing long-term dynamics, as the constant $C$ involved is not uniformly bounded over time. In many applications, our interest extends beyond approximating the transient behavior of an SDE to understanding the long-term statistical properties of its numerical discretization. Firstly, we aim to demonstrate that the discretized scheme is ergodic with respect to some invariant measure, which typically depends on the timestep $\Delta t$. Even when the continuous-time system has an invariant measure, the discretized version may not retain this property.

\begin{example}
Consider the SDE
\[
dX_t = -X(t)^3\,dt + dW_t,
\]
which is ergodic with invariant measure
\[
\pi(x) \propto \exp\left(-\frac{1}{2}x^4\right).
\]
Applying the Euler-Maruyama discretization yields
\[
X_{n+1} = X_n - \Delta t X_n^3 + \Delta W_n.
\]
However, there remains a small but non-negligible probability that the noise term $\Delta W_n$ can push the solution into a regime where $X_n \approx \Delta t^{-1}$, leading to explosive behavior and thus loss of ergodicity for the numerical scheme.
\end{example}
Much like in the continuous case, discussed in section \ref{sec: linfty}, the convergence to a stationary distribution can often be established via two key assumptions: a Lyapunov condition (\ref{ass:lyapunov}) and a minorization condition (\ref{ass:minorization}). In practice, establishing the minorization condition is often relatively straightforward, whereas verifying the Lyapunov condition requires assumptions on the drift term\cite{mattingly2002ergodicity}. A common choice for the Lyapunov function is $\mathcal{K}(x) = 1 + |x|^n$ for some $n > 0$. Following this, we can apply theorem \ref{thm:exp-conv-linf} to prove exponential convergence of the iterate distribution to the invariant distribution. Once ergodicity and irreducibility of the discretized scheme have been established, we can invoke results (e.g., \cite[Chapter~17]{Meyn2012-oq}) to obtain almost sure convergence:
\begin{equation}
\widehat{\varphi}_{N_{\text {iter }}, \Delta t} \xrightarrow[N_{\text {iter }} \rightarrow+\infty]{\text { a.s. }} \int_{\mathcal{X}} \varphi \mathrm{d} \pi_{\Delta t} .
\end{equation}
where $\pi_{\Delta t}$ denotes the invariant measure associated with the discretized process.
%\[
%\widehat{\varphi}_{N_{\text{iter}},\Delta t} \to \int_{\chi} \varphi \, d\pi_{\Delta t} \quad %\text{as} \quad N_{\text{iter}} \to +\infty,
%\]
Suppose now, that Theorem \ref{thm:exp-conv-linf} holds for some ergodic measure $\pi_{\Delta t}$. Assuming $x^0 \sim \pi_{\Delta t}$ (and so $x^n \sim \pi_{\Delta t}$ for all $n\geq 1$) we can use the stationarity property
\begin{equation}
\mathbb{E}_{\pi_{\Delta t}}\left[\Pi_{\Delta t} \varphi\left(x^n\right) \Pi_{\Delta t} \varphi\left(x^m\right)\right]=\mathbb{E}_{\pi_{\Delta t}}\left[\Pi_{\Delta t} \varphi\left(x^{|n-m|}\right) \Pi_{\Delta t} \varphi\left(x^0\right)\right]
\end{equation}
to derive expressions for the asymptotic variance of the discretization, similar to that of the continuous SDE above:
\begin{equation}
\varsigma_{\varphi, \Delta t}^2=N_{\mathrm{corr}, \Delta t, \varphi} \varsigma_{\varphi, \mathrm{iid}, \Delta t}^2 \quad \text { where } \varsigma_{\varphi, \mathrm{iid}, \Delta t}^2=\int_{\mathcal{X}}\left(\Pi_{\Delta t} \varphi\right)^2 \mathrm{~d} \pi_{\Delta t},
\end{equation}
where $\Pi_{\Delta t} \varphi=\varphi-\int_{\mathcal{X}} \varphi \mathrm{d} \pi_{\Delta t}$, $\varsigma_{\varphi, \mathrm{iid}, \Delta t}^2$ refers to the asymptotic variance of i.i.d. samples and $N_{\mathrm{corr}, \Delta t, \varphi}$ quantifies the number of steps required for the process to decorrelate from its previous state. Furthermore, in the limit of small $\Delta t$, the relation between the decorrelation of discrete and continuous dynamics is
$$\theta_{\mathrm{corr}, \varphi} \simeq N_{\mathrm{corr}, \Delta t, \varphi} \Delta t.$$



\begin{comment}
    \begin{assumption}[Lyapunov Condition] \label{lyapass}
There exists a function $\mathcal{K}: \chi \to [1,+\infty)$ and a constant $\alpha \in (0,1)$ such that
\[
(P\mathcal{K})(x) \leq \alpha \mathcal{K}(x) + R, \quad \forall x \in \chi,
\]
for some $R > 0$.
\end{assumption}

\begin{assumption}[Minorization Condition] \label{minass}
There exists a constant $\eta \in (0,1)$ and a probability measure $\lambda$ such that
\[
\inf_{x \in \mathcal{C}} P(x, dy) \geq \eta \lambda(dy),
\]
where the set $\mathcal{C}$ is defined by
\[
\mathcal{C} = \left\{x \in \mathcal{X} \,\middle|\, \mathcal{K}(x) \leq \mathcal{K}_{\max}\right\},
\]
for some $\mathcal{K}_{\max} > 1 + \frac{2R}{1-\alpha}$, with $\alpha$ and $R$ as introduced in Assumption~\ref{lyapass}.
\end{assumption}
\end{comment}




%\subsubsection{Exponential Convergence in Weighted $L^\infty$ Spaces}
\begin{comment}
    Define the weighted $L^\infty$ norm by
\[
\|\varphi\|_{L^\infty_\mathcal{K}} = \left\| \frac{\varphi}{\mathcal{K}} \right\|_{L^\infty}.
\]
The space $L^\infty_\mathcal{K}(\chi)$ consists of measurable functions for which this norm is finite.

\begin{theorem}[Exponential Convergence in Weighted $L^\infty$ Spaces] \label{expoconvthm}
Suppose Assumptions~\ref{lyapass} and~\ref{minass} hold. Then, the Markov kernel $P$ admits a unique invariant probability measure $\pi$, satisfying
\[
\int_{\chi} \mathcal{K}(x) \, d\pi(x) < +\infty.
\]
Moreover, there exist constants $C > 0$ and $r \in (0,1)$ such that, for any $\varphi \in L^\infty_\mathcal{K}(\chi)$ and any $n \in \mathbb{N}$,
\[
\left\| P^n \varphi - \int_{\chi} \varphi \, d\pi \right\|_{L^\infty_\mathcal{K}} \leq C r^n \left\| \varphi - \int_{\chi} \varphi \, d\pi \right\|_{L^\infty_\mathcal{K}}.
\]
\end{theorem}

\subsubsection{Almost Sure Convergence of Time Averages}

Given a test function $\varphi$, we define the time average computed along the discretized trajectory as
\[
\widehat{\varphi}_{N_{\text{iter}},\Delta t} = \frac{1}{N_{\text{iter}}} \sum_{n=0}^{N_{\text{iter}}-1} \varphi(x^n).
\]
\end{comment}


\subsection{Error Analysis: Decomposition}

In practice, we are interested in using $\widehat{\varphi}_{N_{\text{iter}},\Delta t}$ to approximate
\[
\mathbb{E}_\pi(\varphi) = \int_\chi \varphi \, d\pi.
\]
The total error can be decomposed into two contributions:
\begin{equation} \label{eq:statistic_bias}
    \widehat{\varphi}_{N_{\text{iter}},\Delta t} - \mathbb{E}_\pi(\varphi) = \underbrace{\left( \widehat{\varphi}_{N_{\text{iter}},\Delta t} - \mathbb{E}_{\pi_{\Delta t}}(\varphi) \right)}_{\text{Statistical Error}} + \underbrace{\left( \mathbb{E}_{\pi_{\Delta t}}(\varphi) - \mathbb{E}_\pi(\varphi) \right)}_{\text{Bias}}.
\end{equation}

\paragraph{Statistical Error.}
Once the asymptotic regime is attained, the statistical error scales like $\mathcal{O}(1/\sqrt{N_{\text{iter}}})$. More precisely, for asymptotic variance $\varsigma_{\varphi,\Delta t}^2$
\[
\lim_{\Delta t \to 0} \Delta t \, \varsigma_{\varphi,\Delta t}^2 = \varsigma_\varphi^2,
\]
where
\[
\varsigma_\varphi^2 = \lim_{t \to +\infty} t \, \mathrm{Var}_\pi(\hat{\varphi}_t), \quad \varsigma_{\varphi,\Delta t}^2 = \lim_{N_{\text{iter}} \to +\infty} N_{\text{iter}} \, \mathrm{Var}_{\pi_{\Delta t}}\left( \widehat{\varphi}_{N_{\text{iter}},\Delta t} \right).
\]

\paragraph{Bias (Systematic Error).}
The bias arises because $\pi_{\Delta t} \neq \pi$. It can be expanded asymptotically as
\begin{equation} \label{bias_error}
    \int_{\chi} \varphi \, d\pi_{\Delta t} = \int_{\chi} \varphi \, d\pi + \Delta t^p \int_{\chi} \varphi f \, d\pi + \mathcal{O}(\Delta t^{p+1}),
\end{equation}
for some function $f$ and some $p \geq \alpha$, where $\alpha$ is the weak order of the numerical scheme. The bias becomes noticeable only when the total simulation time $T = N_{\text{iter}} \Delta t$ is large.
\subsection{Analysis of the Bias}

We now provide a detailed justification of the bias expansion (\ref{bias_error}) under suitable regularity and stability assumptions. This expansion captures the leading-order contribution of the difference between the invariant measure of the numerical scheme $\pi_{\Delta t}$ and the true invariant measure $\pi$. We assume that the transition operator $P_{\Delta t}$ of the numerical scheme admits an asymptotic expansion in powers of $\Delta t$:
\[
P_{\Delta t} \varphi = \varphi + \Delta t \mathcal{A}_1 \varphi + \Delta t^2 \mathcal{A}_2 \varphi + \cdots + \Delta t^{p+1} \mathcal{A}_{p+1} \varphi + \Delta t^{p+2} r_{\varphi, \Delta t},
\]
for test functions $\varphi$ in a suitable space $\mathcal{S}$. The operators $\mathcal{A}_k$ are typically differential operators determined by the structure of the numerical scheme, and the order $p$ corresponds to the weak accuracy of the integrator. Typically:
\begin{itemize}
    \item $\mathcal{A}_k$ are identified by expanding $\varphi(x^{n+1})$ around $\varphi(x^n)$ using the dynamics of the scheme.
    \item $\mathcal{A}_1 = \mathcal{L}$, where $\mathcal{L}$ is the infinitesimal generator of the continuous-time Markov process.
    \item More generally, if the numerical scheme has weak order $p$, then $\mathcal{A}_k = \frac{\mathcal{L}^k}{k!}$ for $1 \leq k \leq p$.
\end{itemize}  
We work in a class of smooth test functions with controlled growth.
\begin{definition}[Space of Smooth Functions]
Let $(\mathcal{K}_n)_{n \in \mathbb{N}}$ be a sequence of weight functions $\mathcal{K}_n : \mathcal{X} \to [1, +\infty)$, such that
\[\mathcal{K}_n \leq \mathcal{K}_{n+1}, \forall n \geq 0.\]
Define:
\[
\mathcal{S} = \left\{ \varphi \in C^\infty(\mathcal{X}) \,\middle|\, \forall k \in \mathbb{N}^d, \exists m \in \mathbb{N} \text{ such that } \partial^k \varphi \in L^\infty_{\mathcal{K}_m}(\mathcal{X}) \right\},
\]
where $L^\infty_{\mathcal{K}_m}$ denotes the space of functions satisfying $\|\partial^k \varphi / \mathcal{K}_m\|_{L^\infty} < \infty$.

\begin{itemize}
    \item For bounded $\mathcal{X}$, choose $\mathcal{K}_n = 1$, so $\mathcal{S} = C^\infty(\mathcal{X})$.
    \item For unbounded domains, a typical choice is $\mathcal{K}_n(x) = 1 + |x|^n$, ensuring polynomial growth bounds.
\end{itemize}
\end{definition}

In the case where $\mathcal{X}$ is bounded, we may take $\mathcal{K}_n = 1$ and recover $\mathcal{S} = C^\infty(\mathcal{X})$. For unbounded domains, a typical choice is $\mathcal{K}_n(x) = 1 + |x|^n$, ensuring that functions and their derivatives grow at most polynomially. To construct the bias expansion, we require control over solutions of equations involving the operator $\mathcal{A}_1$ and its adjoint.
\begin{assumption}[Stability of Smooth Functions by Inverse Operators] \label{ass:inverse_stability}
Define the subspace of zero-mean test functions:
\[
\mathcal{S}_0 = \left\{ \phi \in \mathcal{S} \,\middle|\, \int_\chi \phi \, d\pi = 0 \right\}.
\]
Assume that $\mathcal{S}$ is dense in $L^2(\pi)$ (in particular, $\mathcal{K}_n \in L^2(\pi)$ for all $n$), and that the inverse operators $\mathcal{A}_1^{-1}$ and $(\mathcal{A}_1^*)^{-1}$ preserve $\mathcal{S}_0$.
\end{assumption}
That is, for $\phi \in \mathcal{S}_0$, the unique solutions $\Phi$ to $\mathcal{A}_1 \Phi = \phi$ and $\mathcal{A}_1^* \Phi = \phi$ also belong to $\mathcal{S}_0$. As mentioned, typically, $\mathcal{A}_1 = \mathcal{L}$. We now state a rigorous result quantifying the bias between the invariant measures $\pi_{\Delta t}$ of the numerical scheme and the target invariant measure $\pi$ of the continuous-time dynamics.

\begin{theorem}[Error estimates on the invariant measure, \texorpdfstring{\cite[Theorem~3.3]{Lelievre2016-hx}}{}]
    \label{thm:error-inv-measure}
    Suppose that the stability assumption \ref{ass:inverse_stability} is satisfied, and that the numerical transition operator $P_{\Delta t}$ admits an expansion
    \begin{equation} \label{A_expansion}
        P_{\Delta t} \varphi = \varphi + \Delta t \mathcal{A}_1 \varphi + \cdots + \Delta t^{p+1} \mathcal{A}_{p+1} \varphi + \Delta t^{p+2} r_{\varphi,\Delta t},
    \end{equation}
    for any $\varphi \in \mathcal{S}$, with some $p \in \mathbb{N}$. Assume that there exist constants $K > 0$, $m \in \mathbb{N}$, and $\Delta t^* > 0$ such that the remainder is uniformly bounded:
    \[
    \| r_{\varphi,\Delta t} \|_{L^\infty(\mathcal{K}_m)} \leq K, \quad \text{for all } \Delta t \leq \Delta t^*.
    \]
    Additionally, assume the following:
    \begin{enumerate}
        \item For each $k \geq 1$, the operator $\mathcal{A}_k$ leaves $\mathcal{S}$ invariant;
        \item For all $k \in \{1, \ldots, p\}$, the operator $\mathcal{A}_k$ satisfies
        \begin{equation} \label{Ak_mean_zero}
            \int_\mathcal{X} \mathcal{A}_k \varphi \, d\pi = 0 \quad \text{for all } \varphi \in \mathcal{S};
        \end{equation}
        \item The function $g_{p+1} := \mathcal{A}_{p+1}^* 1 \in \mathcal{S}_0$;
        \item The measure $\pi_{\Delta t}$ integrates all scale functions: for all $n \in \mathbb{N}$,
        \[
        \int_\mathcal{X} \mathcal{K}_n \, d\pi_{\Delta t} < \infty.
        \]
    \end{enumerate}
    Then, there exists $L > 0$ such that, for all $\Delta t \in (0, \Delta t^*]$,
    \begin{equation} \label{bias_error_estimate}
        \int_\mathcal{X} \varphi \, d\pi_{\Delta t} = \int_\mathcal{X} \varphi \, d\pi + \Delta t^p \int_\mathcal{X} \varphi f_{p+1} \, d\pi + \Delta t^{p+1} R_{\varphi,\Delta t},
    \end{equation}
    with $|R_{\varphi,\Delta t}| \leq L$, and where the bias function $f_{p+1} \in \mathcal{S}_0$ is given by
    \[
    f_{p+1} = -(\mathcal{A}_1^*)^{-1} g_{p+1}.
    \]
    That is,
    \[
    \mathcal{A}_1^* f_{p+1} = -\mathcal{A}_{p+1}^* 1.
    \]
\end{theorem}
\begin{proof} [Outline of proof]
    The proof is in two parts. First, the result is proven for functions $\varphi \in (P_{\Delta t} - \Id)\mathcal{S}.$ Afterwards, the proof is extended to all functions in $\mathcal{S}_0$. Consider $\phi \in \mathcal{S}$. Then invariance of $\pi_{\Delta t}$ by the discretized dynamics gives
    \[
    \int_{\mathcal{X}} P_{\Delta t} \phi \, d\pi_{\Delta t} = \int_{\mathcal{X}} \phi \, d\pi_{\Delta t},
    \]
    from which it follows, that
    $$
        \int_\chi\left[\left(\frac{P_{\Delta t}-\mathrm{Id}}{\Delta t}\right) \phi\right] d \pi_{\Delta t} = 0
    $$
    Then, by (\ref{Ak_mean_zero}) and using the fact, that
    \[
    \frac{P_{\Delta t} - \mathrm{Id}}{\Delta t} = \mathcal{A}_1 + \Delta t \mathcal{A}_2 + \Delta t^2 \mathcal{A}_3 + \cdots,
    \]
    we can expand
    $$
    \begin{aligned}
        \int_{\mathcal{X}} & {\left[\left(\frac{P_{\Delta t}-\mathrm{Id}}{\Delta t}\right) \phi\right]\left(1+\Delta t^p f\right) \mathrm{d} \pi } \\
        & =\Delta t^p \int_{\mathcal{X}}\left(\mathcal{A}_{p+1} \phi+\left(\mathcal{A}_1 \phi\right) f\right) \mathrm{d} \pi \\
        & \quad+\Delta t^{p+1} \int_{\mathcal{X}}\left(r_{\phi, \Delta t}+\left[\frac{P_{\Delta t}-\mathrm{Id}-\Delta t \mathcal{A}_1}{\Delta t^2} \phi\right] f\right) \mathrm{d} \pi .
    \end{aligned}
    $$
    The first integrand on the right hand side forces
    \[f = f_{p+1} = -(\mathcal{A}_1^*)^{-1} g_{p+1}\]
    and so (\ref{bias_error_estimate}) holds for $\varphi = (\Id - P_{\Delta t})\phi / \Delta t$. The next step is to prove for all functions in $\mathcal{S}_0$. With that in mind, a natural approach would be to set $\phi = \Delta t (\Id - P_{\Delta t})^{-1} \varphi$. However, the inverse $(\Id - P_{\Delta t})$ is well-defined only on spaces with mean 0 with respect to $\pi_{\Delta t}$. When it is well-defined, there is no control over its derivatives.
    \begin{mdframed}[hidealllines=true,backgroundcolor=yellow!20]
    Finish this - either remove proof or explain well.
    \end{mdframed}
\end{proof}


\bibliographystyle{plain}
\bibliography{references}

\end{document}
